\documentclass[11pt,a4paper]{amsart}
\usepackage[utf8]{inputenc}
\usepackage{amsmath,amssymb,amsthm,mathtools}
\usepackage[margin=1.1in]{geometry}
\usepackage{booktabs}
\usepackage[colorlinks=true,linkcolor=blue,citecolor=blue,urlcolor=blue]{hyperref}
\usepackage{xurl}

\theoremstyle{plain}
\newtheorem{theorem}{Theorem}[section]
\newtheorem{lemma}[theorem]{Lemma}
\newtheorem{proposition}[theorem]{Proposition}
\newtheorem{corollary}[theorem]{Corollary}
\theoremstyle{definition}
\newtheorem{definition}[theorem]{Definition}
\newtheorem{example}[theorem]{Example}
\newtheorem{remark}[theorem]{Remark}

\newcommand{\A}{\mathcal{A}}
\newcommand{\B}{\mathcal{B}}
\newcommand{\Cc}{\mathcal{C}}
\newcommand{\Mm}{\mathcal{M}}
\newcommand{\Nn}{\mathcal{N}}
\newcommand{\Yy}{\mathcal{Y}}
\newcommand{\Avg}{\operatorname{Avg}}
\newcommand{\Cp}{\Cc^{+}}
\DeclareMathOperator{\wt}{w}

\begin{document}

\title[Averaging in separating union-closed families of height four]
{Averaging in separating union-closed families of height four:\\ the height-four case}

\author{Yavuz Selim K\i l\i n\c{c}}
\thanks{Preprint and verification package archived at Zenodo:
\url{https://doi.org/10.5281/zenodo.21739198}.}
\address{Independent researcher, T\"urkiye}
\email{yavuzselim18mart@gmail.com}

\subjclass[2020]{05D05, 06A07}
\keywords{Frankl's conjecture, union-closed sets conjecture, averaging, separating family, height}

\begin{abstract}
Let $\A\subseteq 2^{[n]}$ be a separating union-closed family whose longest chain
consists of exactly four member sets. Bouchard proved that
$\Avg(\A)\ge n/2$ holds for such families when $|\B|\le 2$, where $\B$ is a smallest
irredundant subfamily of $\A_{<n/2}$ spanning the same base set, and left the
range $3\le|\B|\le4$ open. We prove the full height-four case in its natural
range $4\le n$.

Writing $T$ for the set of $x$ with $[n]\setminus\{x\}\in\A$ and $t=|T|$, we first
prove structural lemmas about the core
$\Cc=\A\setminus(\{[n]\}\cup\{[n]\setminus\{x\}:x\in T\})$ and deduce a master
inequality
\[
  \sum_{\substack{A\in\A\\2|A|<n}}\bigl(n-2|A|\bigr)
  \le
  (t+1)\,\bigl|\{A\in\A:2|A|<n\}\bigr|.
\]
This gives $\Avg(\A)\ge n/2$ whenever $n\ge4t+4$, and reduces the remaining
problem to the window $n\le4t+3$. We then introduce a pair-line structure on the
small member sets $\Cp=\{A\in\A:2|A|<n\}$: two small sets either lie on a unique
maximal core line or form a rigid external pair. Combining line capacity, a
Fisher-type count of maximal core complements, and a local bound around a
largest small set, we prove the residual inequality
\[
  \sum_{A\in\Cp}(n-2|A|)\le t(n-2)+n
\]
for all admissible $n,t$ with $4\le n$. Consequently
\[
  \Avg(\A)\ge \frac n2
\]
for every separating union-closed family of height four with $4\le n$.
In particular, Frankl's conjecture holds for this class.
\end{abstract}

\maketitle

\section{Introduction}

A finite family $\A$ of finite sets is \emph{union-closed} if $X,Y\in\A$ implies
$X\cup Y\in\A$. We always denote its base set $\bigcup_{A\in\A}A$ by
$[n]=\{1,\dots,n\}$ and assume $|\A|\ge2$; union-closedness then gives $[n]\in\A$.
Frankl's union-closed sets conjecture~\cite{Frankl95,BruhnSchaudt} asserts that
some element of $[n]$ lies in at least $|\A|/2$ member sets. Despite the constant
lower bounds obtained after Gilmer's breakthrough~\cite{Gilmer,Sawin,AHS}, the
conjecture itself remains open.

Following the standard double-counting identity
$\sum_{x\in[n]}|\A_x|=\sum_{A\in\A}|A|$ (see~\cite[\S6]{BruhnSchaudt}), the
\emph{averaging approach} replaces the conjecture by the stronger statement
\begin{equation}\label{eq:avg}
  \Avg(\A):=\frac1{|\A|}\sum_{A\in\A}|A|\ge\frac n2.
\end{equation}
Inequality~\eqref{eq:avg} is false for general union-closed families, and
Reimer's average set size theorem~\cite{Reimer} shows what can be salvaged in
general; the scope of averaging has been mapped out by Cz\'edli and
coauthors~\cite{Czedli09,CzedliMS} and, recently, in a formally verified
setting~\cite{HachimoriKashiwabara}. The natural class on which~\eqref{eq:avg}
does have a chance is that of \emph{separating} families: $\A$ is separating if
for all distinct $x,y\in[n]$ some $A\in\A$ contains exactly one of them.
Restricting the conjecture to separating families costs nothing.

A \emph{chain} in $\A$ is a subfamily that is totally ordered by inclusion, and
the \emph{height} $h=h(\A)$ is the maximum number of member sets in a chain.
(We follow the convention of~\cite{Bouchard,Tian}: the height counts \emph{sets},
not covering relations. Thus $h\le3$ means that no four member sets form a
chain.) Frankl's conjecture is known for $h\le3$~\cite{Tian}, and for separating
families the stronger inequality~\eqref{eq:avg} is known in that range as
well~\cite[Thm.~1.4]{Bouchard}. The height-four case is the first genuinely
nontrivial one, and it is also the last one for which averaging can possibly
work: Bouchard constructs, for every $h\ge5$ and every sufficiently large $n$,
separating union-closed families with $\Avg(\A)<n/2$~\cite[Thm.~3.2]{Bouchard}.
In other words, $h=4$ is exactly the boundary of the method.

For $h=4$ the previous state of the art was~\cite[Thm.~2.1]{Bouchard}: writing
$\B$ for a smallest irredundant subfamily of $\A_{<n/2}$ with the same base set
as $\A_{<n/2}$, one always has $0\le|\B|\le h$, and~\eqref{eq:avg} holds for
separating union-closed families with $h=4\le n$ and $|\B|\le2$. The remaining
cases $|\B|\in\{3,4\}$ were open. The present paper closes them.

\subsection*{Results of this paper}

We attack the class through a different parameter. Let
\[
  T:=\{x\in[n]:[n]\setminus\{x\}\in\A\},\qquad t:=|T|,
\]
\[
  Z:=[n]\setminus T,\qquad m:=|Z|=n-t,
\]
let $\Yy:=\{[n]\setminus\{x\}:x\in T\}$ be the family of \emph{coatoms}, and let
\[
  \Cc:=\A\setminus(\{[n]\}\cup\Yy)
\]
be the \emph{core}. Section~\ref{sec:struct} proves structural facts about $\Cc$
(Lemmas~\ref{lem:under}--\ref{lem:P}), of which the key new ingredient is that
every maximal member of the core misses at most one element of $Z$
(Lemma~\ref{lem:Zminus}). These combine into a pair inequality
(Lemma~\ref{lem:P}) and yield the master inequality
(Theorem~\ref{thm:master}).

For separating union-closed families of height four with $4\le n$, the master
inequality gives several consequences: $\Avg(\A)\ge n/2$ whenever
$n\ge4t+4$ (Theorem~\ref{thm:A}), whenever
$(t+1)q\le t(n-2)+n$ (Theorem~\ref{thm:B}), and whenever $t\le2$
(Theorem~\ref{thm:t2-final}). The large-$n$ theorem confines the problem to the
window $n\le4t+3$.

The main new part is Section~\ref{sec:closing}. There we analyse the pair
structure of
\[
  \Cp:=\{A\in\A:2|A|<n\}.
\]
Two distinct members of $\Cp$ either lie on a unique maximal core line or form a
rigid external pair (Lemma~\ref{lem:pair}). This line structure, together with a
capacity bound (Lemma~\ref{lem:cap}), a Fisher-type count
(Lemmas~\ref{lem:fisher}--\ref{lem:q}), and a local bound around a largest small
set (Lemma~\ref{lem:local}), proves the residual inequality
\[
  \Phi^+:=\sum_{A\in\Cp}(n-2|A|)\le t(n-2)+n
\]
for all admissible $n,t$ with $4\le n$ (Theorem~\ref{thm:window} and
Corollary~\ref{cor:residual}). This residual inequality implies
$\Avg(\A)\ge n/2$, since $\Phi\le\Phi^+$ and
Proposition~\ref{prop:reform} characterises the average bound by
$\Phi\le t(n-2)+n$.

Thus the main theorem is the following.

\begin{theorem}[Main theorem; Theorem~\ref{thm:main}]
Let $\A\subseteq2^{[n]}$ be separating and union-closed with $h(\A)=4\le n$. Then
\[
  \Avg(\A)\ge\frac n2.
\]
In particular, Frankl's conjecture holds for every separating union-closed
family of height four with $4\le n$.
\end{theorem}

Two features of the proof are worth noting. First, it is insensitive to
$|\B|$: the extremal family of Example~\ref{ex:sharp} has $|\B|=3$, i.e.\ it lies
in the range left open in~\cite{Bouchard}, and the master inequality covers it
with equality. Second, the simple counting bound $q\le n-2$ is not available in
the residual window: Proposition~\ref{prop:qcex} constructs, for every even
$n\ge6$, a separating union-closed family of height four with $t=n$, $q=n$, and
$\Phi^+=2n$. The proof therefore bounds $\Phi^+$ directly, line by line, rather
than through $q$ alone.

\section{Preliminaries}\label{sec:prelim}

Throughout, $\A\subseteq2^{[n]}$ is union-closed with base set $[n]$ and
$|\A|\ge2$, so $[n]\in\A$. For $A\subseteq[n]$ we write $A^c:=[n]\setminus A$
and define the \emph{weight}
\[
  \wt(A):=n-2|A|.
\]
Thus $\wt(A)>0$ exactly when $|A|<n/2$. We abbreviate
\[
  \Phi:=\sum_{A\in\Cc}\wt(A),\qquad
  \Cp:=\{A\in\Cc:\wt(A)>0\},\qquad
  q:=|\Cp|,\qquad
  \Phi^+:=\sum_{A\in\Cp}\wt(A).
\]
Note that $\Cp=\{A\in\A:2|A|<n\}=\A_{<n/2}$ whenever $n\ge3$, because $[n]$ and
the coatoms have size at least $n-1\ge n/2$. Clearly $\Phi\le\Phi^+$.

\begin{proposition}\label{prop:reform}
$\Avg(\A)\ge n/2$ if and only if $\Phi\le t(n-2)+n$.
\end{proposition}

\begin{proof}
$\Avg(\A)\ge n/2$ is equivalent to $\sum_{A\in\A}(n-2|A|)\le0$. Splitting the sum
over $\{[n]\}$, $\Yy$ and $\Cc$ gives
\[
  \sum_{A\in\A}\wt(A)
  =(n-2n)+t\bigl(n-2(n-1)\bigr)+\Phi
  =\Phi-n-t(n-2). \qedhere
\]
\end{proof}

We record the basic fact underlying everything, due to Bouchard
\cite[Lem.~1.3]{Bouchard}; we include the short proof to keep the paper
self-contained.

\begin{lemma}\label{lem:coatom}
Let $\A$ be separating and union-closed with $|\A|\ge2$. Then every maximal chain
of $\A$ contains a member set of size $n-1$. In particular $t\ge1$.
\end{lemma}

\begin{proof}
Let $[n]=K_1\supsetneq K_2\supsetneq\cdots$ be a maximal chain and suppose
$|K_2|\le n-2$. Choose distinct $x,y\in K_2^c$. Since $\A$ is separating there
is $A\in\A$ containing exactly one of them, say $x\in A$ and $y\notin A$. Then
$K_2\subsetneq K_2\cup A$ because $x\notin K_2$, and $K_2\cup A\subsetneq[n]$
because $y\notin K_2\cup A$. Inserting $K_2\cup A\in\A$ into the chain
contradicts maximality.
\end{proof}

\section{Structure of the core}\label{sec:struct}

Throughout this section $\A$ is separating, union-closed, and $h(\A)=4$. Let
$\Mm$ denote the set of maximal members of $\Cc$ with respect to inclusion and
$\Nn:=\Cc\setminus\Mm$.

\begin{lemma}\label{lem:under}
Every $A\in\Cc$ satisfies $A\subsetneq Y$ for some coatom $Y\in\Yy$. Equivalently,
$T\setminus A\neq\emptyset$ for every $A\in\Cc$.
\end{lemma}

\begin{proof}
Extend $\{A,[n]\}$ to a maximal chain. By Lemma~\ref{lem:coatom} that chain
contains a set $Y$ with $|Y|=n-1$, and $Y\in\Yy$ by definition of $T$. Since
$A\in\Cc$ we have $|A|\le n-2$, so $A$ and $Y$ are comparable with
$A\subsetneq Y$. Writing $Y=[n]\setminus\{x\}$ with $x\in T$ gives
$x\in T\setminus A$.
\end{proof}

\begin{lemma}\label{lem:nochain}
$\Cc$ contains no chain of three sets. Consequently every member of $\Cc$ is
either maximal or minimal in $\Cc$, and every $N\in\Nn$ satisfies
$N\subsetneq M$ for some $M\in\Mm$.
\end{lemma}

\begin{proof}
Suppose $A\subsetneq B\subsetneq C$ with $A,B,C\in\Cc$. By Lemma~\ref{lem:under}
there is a coatom $Y$ with $C\subsetneq Y$. Then
$A\subsetneq B\subsetneq C\subsetneq Y\subsetneq[n]$ is a chain of five member
sets, contradicting $h(\A)=4$. A poset without three-element chains has all
elements maximal or minimal.
\end{proof}

The next lemma is the pivot of the structural part.

\begin{lemma}\label{lem:Zminus}
Every $M\in\Mm$ satisfies $|Z\setminus M|\le1$. In particular $|M|\ge m-1$ and
\[
  \wt(M)\le n-2(m-1)=2t+2-n.
\]
\end{lemma}

\begin{proof}
Suppose $z_1,z_2\in Z\setminus M$ are distinct. Since $\A$ is separating, some
$A\in\A$ contains exactly one of $z_1,z_2$. Now $A\neq[n]$, and $A$ is not a
coatom either, because coatoms are of the form $[n]\setminus\{x\}$ with $x\in T$
and therefore contain all of $Z$. Hence $A\in\Cc$.

If $A\subseteq M$ then $z_1,z_2\notin A$, a contradiction. So
$A\not\subseteq M$, whence $M\subsetneq M\cup A\in\A$. By maximality of $M$ in
$\Cc$ we get $M\cup A\notin\Cc$, i.e.\ $M\cup A\in\Yy\cup\{[n]\}$. Every member
of $\Yy\cup\{[n]\}$ contains $Z$, so $Z\setminus M\subseteq A$ and in particular
$z_1,z_2\in A$ --- again a contradiction.
\end{proof}

\begin{lemma}\label{lem:pairM}
For distinct $M,M'\in\Mm$ we have $M\cup M'\in\Yy\cup\{[n]\}$. Consequently
$M^c\cap M'^c$ is contained in $T$ and has at most one element; in particular,
distinct maximal core sets omit distinct elements of $Z$.
\end{lemma}

\begin{proof}
$M\subsetneq M\cup M'\in\A$, so maximality gives
$M\cup M'\in\Yy\cup\{[n]\}$, i.e.\ $|(M\cup M')^c|\le1$ and, if nonempty,
$(M\cup M')^c\subseteq T$. Now $(M\cup M')^c=M^c\cap M'^c$.
\end{proof}

\begin{lemma}\label{lem:children}
Let $M\in\Mm$ and let $\Nn_M:=\{N\in\Cc:N\subsetneq M\}$. Then:
\begin{enumerate}
\item[(i)] any two distinct $N,N'\in\Nn_M$ satisfy $N\cup N'=M$;
\item[(ii)] the sets $M\setminus N$, $N\in\Nn_M$, are pairwise disjoint and
nonempty, hence
\[
  \sum_{N\in\Nn_M}|M\setminus N|\le|M|
  \qquad\text{and}\qquad
  |\Nn_M|\le|M|.
\]
\end{enumerate}
\end{lemma}

\begin{proof}
(i) $N\cup N'\in\A$ and $N\cup N'\subseteq M$, so $N\cup N'\in\Cc$. If
$N\cup N'=N$ then $N'\subsetneq N\subsetneq M$ is a three-chain in $\Cc$,
contradicting Lemma~\ref{lem:nochain}; symmetrically $N\cup N'\neq N'$. If
$N\cup N'\subsetneq M$ then $N\subsetneq N\cup N'\subsetneq M$ is again a
three-chain. Hence $N\cup N'=M$.

(ii) $N\cup N'=M$ is equivalent to
$(M\setminus N)\cap(M\setminus N')=\emptyset$, and $M\setminus N\neq\emptyset$
since $N\subsetneq M$.
\end{proof}

We can now state the inequality that drives the master argument.

\begin{lemma}[Pair inequality]\label{lem:P}
Any two distinct $A,A'\in\Cc$ satisfy
\[
  |A|+|A'|\ge m-1.
\]
\end{lemma}

\begin{proof}
If $A$ and $A'$ are comparable, say $A\subsetneq A'$, then $A'$ is not minimal
in $\Cc$, hence maximal by Lemma~\ref{lem:nochain}, so $|A'|\ge m-1$ by
Lemma~\ref{lem:Zminus} and the claim follows.

If they are incomparable, then $A\subsetneq A\cup A'\in\A$. If
$A\cup A'\in\Cc$, then $A\cup A'$ is not minimal, hence maximal, so
$|A\cup A'|\ge m-1$ by Lemma~\ref{lem:Zminus}. If $A\cup A'\notin\Cc$, then
$|A\cup A'|\ge n-1\ge m-1$. In both cases
$|A|+|A'|\ge|A\cup A'|\ge m-1$.
\end{proof}

\begin{corollary}\label{cor:empty}
At most one member of $\Cc$ has size smaller than $(m-1)/2$. Moreover, if
$\emptyset\in\A$ then $\Cc=\{\emptyset\}\cup\Mm$, i.e.\ every nonempty core set
is maximal.
\end{corollary}

\begin{proof}
The first claim is immediate from Lemma~\ref{lem:P}. For the second, if
$\emptyset\in\A$ then $\emptyset\in\Cc$, and
$\emptyset\subsetneq N\subsetneq M$ would be a three-chain in $\Cc$ for any
non-maximal nonempty $N$.
\end{proof}

\section{The master inequality}\label{sec:master}

Lemma~\ref{lem:under} guarantees that $T\setminus A\neq\emptyset$ for every
$A\in\Cc$, so the following fractional charge is well defined. For $x\in T$ put
\begin{equation}\label{eq:charge}
  \widetilde C_x
  :=
  \sum_{\substack{A\in\Cp\\ x\in T\setminus A}}
  \frac{\wt(A)}{|T\setminus A|}.
\end{equation}

\begin{proposition}[Exact charge identity]\label{prop:identity}
\[
  \sum_{x\in T}\widetilde C_x=\Phi^+.
\]
\end{proposition}

\begin{proof}
Exchanging the order of summation, each $A\in\Cp$ contributes
$\wt(A)\cdot|T\setminus A|/|T\setminus A|=\wt(A)$.
\end{proof}

\begin{remark}\label{rem:whycharge}
The identity in Proposition~\ref{prop:identity} is exact, with no loss and no
omitted sets. This is the reason for charging $A$ to the elements of
$T\setminus A$ rather than to the elements of $A\cap T$: the latter convention
assigns nothing at all to core sets disjoint from $T$ (for instance to
$\emptyset$, and to every core set when $t=1$, by Lemma~\ref{lem:under}), and
therefore cannot recover $\Phi^+$.
\end{remark}

\begin{theorem}[Master inequality]\label{thm:master}
If $q\ge2$ then
\[
  \Phi^+\le q(t+1).
\]
\end{theorem}

\begin{proof}
Write $\Cp=\{A_1,\dots,A_q\}$ and $a_i:=|A_i|$. Summing Lemma~\ref{lem:P} over
all $\binom q2$ pairs gives
\[
  (q-1)\sum_{i=1}^q a_i
  =
  \sum_{i<j}(a_i+a_j)
  \ge
  \binom q2(m-1),
\]
so
\[
  \sum_{i=1}^q a_i\ge\frac{q(m-1)}2.
\]
Hence
\[
  \Phi^+
  =qn-2\sum_i a_i
  \le qn-q(m-1)
  =q(n-m+1)
  =q(t+1). \qedhere
\]
\end{proof}

\begin{theorem}\label{thm:B}
Assume $4\le n$. If $(t+1)q\le t(n-2)+n$, then $\Avg(\A)\ge n/2$.
\end{theorem}

\begin{proof}
If $q=0$ then $\Phi\le0\le t(n-2)+n$. If $q=1$ then
$\Phi\le\Phi^+\le n\le t(n-2)+n$ because $t\ge1$ and $n\ge4$ give
$t(n-2)\ge0$. If $q\ge2$ then $\Phi\le\Phi^+\le q(t+1)\le t(n-2)+n$ by
Theorem~\ref{thm:master}. In all cases Proposition~\ref{prop:reform} applies.
\end{proof}

\begin{corollary}\label{cor:count}
$\Avg(\A)\ge n/2$ holds whenever $q\le n-2$, and also whenever $t=1$ and
$q\le n-1$.
\end{corollary}

\begin{proof}
The hypothesis of Theorem~\ref{thm:B} reads $q\le n-\frac{2t}{t+1}$, and
$1\le\frac{2t}{t+1}<2$ for $t\ge1$, with value $1$ exactly when $t=1$.
\end{proof}

\begin{example}\label{ex:sharp}
Let $n=4$ and $\A=(2^{[3]}\setminus\{\emptyset\})\cup\{[4]\}$. Then $\A$ is
union-closed and separating, its longest chain
$\{1\}\subsetneq\{1,2\}\subsetneq\{1,2,3\}\subsetneq[4]$ has four sets, and
$\Avg(\A)=16/8=2=n/2$. Here $T=\{4\}$, $t=1$, $m=3$,
$\Cp=\{\{1\},\{2\},\{3\}\}$, $q=3$, and
\[
  \Phi=\Phi^+=6=q(t+1)=t(n-2)+n.
\]
Thus Theorem~\ref{thm:master}, Theorem~\ref{thm:B} and
Proposition~\ref{prop:reform} are simultaneously tight, and
$\widetilde C_4=6=(n-2)+n/t$. This family has $|\B|=3$, so it lies in the range
left open in~\cite{Bouchard}; Theorem~\ref{thm:B} settles it.
\end{example}

\section{The large-\texorpdfstring{$n$}{n} regime}\label{sec:large}

\begin{lemma}\label{lem:kbound}
Assume $n\ge4t+4$. Then:
\begin{enumerate}
\item[(i)] $\Mm\cap\Cp=\emptyset$; in particular every $A\in\Cp$ is properly
contained in some $M\in\Mm$;
\item[(ii)] for every $M\in\Mm$, at most two members of $\Cp$ are proper subsets
of $M$.
\end{enumerate}
\end{lemma}

\begin{proof}
From $n\ge4t+4$ we get $t\le(n-4)/4$ and therefore, by Lemma~\ref{lem:Zminus},
\begin{equation}\label{eq:Mbig}
  |M|\ge m-1=n-t-1\ge n-\frac{n-4}{4}-1=\frac{3n}{4}
  \qquad\text{for every }M\in\Mm.
\end{equation}
In particular $|M|\ge3n/4>n/2$, so $\wt(M)<0$ and $M\notin\Cp$; combined with
Lemma~\ref{lem:nochain} this proves~(i).

For~(ii), fix $M\in\Mm$ and let $N_1,\dots,N_k$ be the members of $\Cp$ with
$N_i\subsetneq M$. Each satisfies $|N_i|<n/2$, hence
$|M\setminus N_i|>|M|-n/2$. By Lemma~\ref{lem:children}(ii) the sets
$M\setminus N_i$ are pairwise disjoint subsets of $M$, so
\[
  k\Bigl(|M|-\frac n2\Bigr)
  <
  \sum_{i=1}^k|M\setminus N_i|
  \le |M|,
\]
i.e.
\[
  k<\frac{|M|}{|M|-n/2}.
\]
The function $x\mapsto x/(x-n/2)$ is decreasing for $x>n/2$, so by~\eqref{eq:Mbig}
the right-hand side is at most
\[
  \frac{3n/4}{3n/4-n/2}=3.
\]
Hence $k\le2$.
\end{proof}

\begin{theorem}\label{thm:A}
If $n\ge4t+4$, then $\Avg(\A)\ge n/2$. More precisely $\Phi\le0$ unless $q=1$,
in which case $\Phi\le2t+2$.
\end{theorem}

\begin{proof}
If $q=0$ then every term of $\Phi$ is nonpositive and $\Phi\le0$.

Assume $q\ge1$. By Lemma~\ref{lem:kbound}(i) we may choose, for each
$A\in\Cp$, a maximal core set $p(A)\in\Mm$ with $A\subsetneq p(A)$. Let
$P:=p(\Cp)\subseteq\Mm$ be the set of chosen parents. By
Lemma~\ref{lem:kbound}(ii) each $M\in P$ is the parent of at most two members of
$\Cp$, so $|P|\ge q/2$. The sets in $P$ are pairwise distinct members of $\Cc$
and, by Lemma~\ref{lem:kbound}(i), none of them lies in $\Cp$; hence
\begin{equation}\label{eq:split}
  \Phi
  \le
  \sum_{A\in\Cp}\wt(A)+\sum_{M\in P}\wt(M)
  =
  \Phi^++\sum_{M\in P}\wt(M),
\end{equation}
because all remaining terms of $\Phi$ are nonpositive. By Lemma~\ref{lem:Zminus}
we have $\wt(M)\le2t+2-n<0$ for every $M\in P$, so
\[
  \sum_{M\in P}\wt(M)\le |P|(2t+2-n)\le\frac q2(2t+2-n).
\]

If $q=1$, then $|P|=1$ and~\eqref{eq:split} gives
$\Phi\le n+(2t+2-n)=2t+2$. Since $n\ge4$ and $t\ge1$ we have
$t(n-2)+n-(2t+2)=t(n-4)+n-2>0$, so $\Phi\le t(n-2)+n$.

If $q\ge2$, then Theorem~\ref{thm:master} applies and~\eqref{eq:split} gives
\[
  \Phi
  \le
  q(t+1)+\frac q2(2t+2-n)
  =
  q\Bigl(2t+2-\frac n2\Bigr)
  \le0,
\]
the last step because $n\ge4t+4$. In all cases $\Phi\le t(n-2)+n$, and
Proposition~\ref{prop:reform} applies.
\end{proof}

\begin{remark}
The proof of Theorem~\ref{thm:A} shows more: for $n\ge4t+4$ the core weight is
already nonpositive, so the coatoms are not needed at all. The threshold is of
the correct shape: the bound $k\le2$ of Lemma~\ref{lem:kbound}(ii) degrades as
$n/t$ decreases, which is precisely the mechanism that makes small $n/t$ the hard
regime.
\end{remark}

\section{Small ground sets and the residual inequality}\label{sec:conj}

By Theorem~\ref{thm:A} we may assume $n\le4t+3$. We first record the exhaustive
verification for small ground sets, then state the residual inequality that is
proved in Sections~\ref{sec:closing}--\ref{sec:main}.

\subsection*{Small ground sets}

The residual inequality can be settled for small $n$ by exhaustive computation.
We carried out two independent computations, described below; both are available
in the repository cited at the end of this section.

\emph{Computation I (full enumeration).} For a fixed $n$ and a canonical
$t$-subset $T\subseteq[n]$ we enumerate every union-closed family $\A$ whose
coatom index set is $T$, pruning only with the definitional constraints (core
sets of size at most $n-2$, union-closedness, no chain of five member sets), so
that the lemmas of Section~\ref{sec:struct} are tested and not presupposed.
Since every quantity involved is invariant under permutations of $[n]$, fixing
$T$ loses nothing. Each family found is checked against
Lemmas~\ref{lem:under}--\ref{lem:P}, Proposition~\ref{prop:identity},
Theorems~\ref{thm:master} and~\ref{thm:A}, the residual inequality
$\Phi^+\le t(n-2)+n$, and the conclusion $\Avg(\A)\ge n/2$. The enumeration is
complete for $n\le5$ and for $n=6$ with $t\le5$; in the labelling of $[n]$ this
covers
\[
  2{,}034\ (n=4),\qquad
  346{,}028\ (n=5),\qquad
  106{,}040{,}636\ (n=6,\ t\le5)
\]
families, with no violation of any statement. For $n=4$ the count $2{,}034$ was
reproduced by an independent brute-force scan of all $2^{16}$ subfamilies of
$2^{[4]}$. The number of families grows too fast for this method to reach
$n\ge7$.

\emph{Computation II (small parts).} Since $\Phi^+$ depends only on the member
sets of size less than $n/2$, one may enumerate the possible small parts $\Cp$
instead of the families themselves. Every family's small part occurs among the
enumerated states, so an upper bound proved this way holds for all families.
A state is discarded only when it is impossible: because the union-closure of
the selected sets would contain a set of size $n-1$ that is not a coatom;
because it would contain a chain of three core sets; because a core set would
contain $T$ (Lemma~\ref{lem:under}); because a set of the closure of size below
$n/2$ was not itself selected; because a closure set properly containing another
one would omit two points of $Z$ (Lemma~\ref{lem:Zminus}); or because two core
sets would violate Lemma~\ref{lem:P}. A branch-and-bound cut discards a subtree
once the accumulated $\Phi^+$ plus all remaining positive weight cannot exceed
$t(n-2)+n$. Adding the pair-structure cuts of Section~\ref{sec:closing} as
further necessary conditions also verifies the cases $n=11$ with $t\le3$.

This gives the following.

\begin{theorem}\label{thm:small}
For the verified range $4\le n$, the inequality
\[
  \Phi^+\le t(n-2)+n
\]
holds for every $4\le n\le10$ and every admissible $t$, and also for $n=11$
with $t\le3$.
Consequently $\Avg(\A)\ge n/2$ for every separating union-closed family of
height four with $4\le n\le10$, and also for the cases
$n=11$, $t\le3$.
\end{theorem}

The restriction $4\le n$ is necessary. On $[3]$, the family
$\{\emptyset,\{1\},\{2\},\{1,2\},[3]\}$ is separating and union-closed,
has height four, and has average size $7/5<3/2$. It lies outside the theorem's
scope; correspondingly its $\Phi^+=5$ exceeds $t(n-2)+n=4$. This example does
not assume that the family is empty-free, and the theorem above makes no such
assumption.

Table~\ref{tab:maxphi} lists the exact maximum of $\Phi^+$ over all admissible
small parts for small $n$, computed with the branch-and-bound cut disabled. On
the fourteen pairs $(n,t)$ for which Computation~I is also complete, the two
computations agree on the maximum in every case, which is a strong consistency
check on the constraint set used in Computation~II.

\begin{table}[ht]
\centering
\begin{tabular}{rrrrrrrrrr}
\toprule
& \multicolumn{9}{c}{$t$}\\
\cmidrule(l){2-10}
$n$ & 1 & 2 & 3 & 4 & 5 & 6 & 7 & 8 & bound at $t=1$\\
\midrule
4 & \textbf{6} & 6 & 6 & 8 & --- & --- & --- & --- & 6\\
5 & 6 & 8 & 8 & 9 & 10 & --- & --- & --- & 8\\
6 & 6 & 8 & 10 & 10 & 10 & 12 & --- & --- & 10\\
7 & 7 & 7 & 11 & 13 & 13 & 13 & 14 & --- & 12\\
8 & 8 & 8 & 10 & 14 & 16 & 16 & 16 & 16 & 14\\
\bottomrule
\end{tabular}
\caption{Exact maximum of $\Phi^+$ over all admissible small parts. The bound
$t(n-2)+n$ is smallest at $t=1$; the only case in which it is attained is
$(n,t)=(4,1)$, realised by Example~\ref{ex:sharp}.}
\label{tab:maxphi}
\end{table}

For $9\le n\le10$ the branch-and-bound run visits between $8\cdot10^6$ and
$4\cdot10^7$ states per value of $t$ and verifies the inequality in every case;
the largest value of $\Phi^+$ encountered there is $22$, against a bound that
ranges from $18$ to $90$. Away from $n=4$ the inequality is therefore very far
from tight.

The code for both computations is available at
\url{https://github.com/YSKpython/frankl-height4-avg-bound}.

\subsection*{The residual inequality}

The residual statement, formerly Conjecture~6.1 in earlier drafts, is the
following.

\begin{theorem}[Residual inequality]\label{thm:residual}
Let $\A\subseteq2^{[n]}$ be separating and union-closed with $h(\A)=4\le n$. Then
\[
  \Phi^+
  =
  \sum_{\substack{A\in\A\\2|A|<n}}(n-2|A|)
  \le
  t(n-2)+n.
\]
\end{theorem}

The proof is given in Sections~\ref{sec:closing} and~\ref{sec:main}: the finite
verification above handles small $n$, Theorem~\ref{thm:A} handles
$n\ge4t+4$, and Section~\ref{sec:closing} closes the remaining window
$n\le4t+3$.

\section{The pair structure of \texorpdfstring{$\Cc^+$}{C+} and the closing argument}
\label{sec:closing}

Throughout this section $\A$ is separating and union-closed with
$h(\A)=4\le n$. We keep the notation of Sections~\ref{sec:prelim}--\ref{sec:struct}
and add the following. Put
\[
  \nu:=\Bigl\lceil\frac n2\Bigr\rceil-1,
  \qquad
  S_A:=[n]\setminus A,
  \qquad
  \lambda_M:=n-|M|=|S_M|.
\]
Thus $A\in\Cp$ if and only if $A\in\Cc$ and $|A|\le\nu$, and
$|S_A|\ge n-\nu=\lfloor n/2\rfloor+1$ for every $A\in\Cp$. For $M\in\Mm$ set
\[
  V_M:=\{A\in\Cp:A\subseteq M\},
  \qquad
  k_M:=|V_M|.
\]
Since $M\in\Cc$ we have $|M|\le n-2$, i.e.\ $\lambda_M\ge2$, and by
Lemma~\ref{lem:Zminus} we have $|M|\ge m-1$, i.e.\ $\lambda_M\le t+1$.

The whole of this section rests on one observation: two member sets of size
below $n/2$ are too large for their union to escape the core. In particular the
argument does not degenerate when $Z$ is small, which is exactly where the
earlier counting arguments lose their grip.

\subsection*{The pair lemma}

\begin{lemma}[Pair lemma]\label{lem:pair}
Let $A\ne A'$ be members of $\Cp$. Then exactly one of the following holds.
\begin{enumerate}
\item[\rm(i)] $A\cup A'\in\Mm$ \hfill \emph{(the pair is internal)};
\item[\rm(ii)] $n$ is odd, $|A|=|A'|=\frac{n-1}{2}$,
$A\cap A'=\emptyset$, and $A\cup A'\in\Yy$ \hfill \emph{(the pair is external)}.
\end{enumerate}
\end{lemma}

\begin{proof}
Since $|A|,|A'|\le\nu$ we have
$|S_A|,|S_{A'}|\ge n-\nu$, whence
\begin{equation}\label{eq:pair}
  |S_A\cap S_{A'}|
  =
  |S_A|+|S_{A'}|-|S_A\cup S_{A'}|
  \ge
  2(n-\nu)-n
  =
  \begin{cases}
    2,& n\text{ even},\\
    1,& n\text{ odd}.
  \end{cases}
\end{equation}
As $|A\cup A'|=n-|S_A\cap S_{A'}|$, two cases arise.

If $|S_A\cap S_{A'}|\ge2$, then $|A\cup A'|\le n-2$, so $A\cup A'$ is neither
$[n]$ nor a coatom, i.e.\ $A\cup A'\in\Cc$. Moreover $A\cup A'$ is not minimal
in $\Cc$: if $A$ and $A'$ are incomparable then
$A\subsetneq A\cup A'$, and if say $A\subsetneq A'$ then
$A\cup A'=A'\supsetneq A$ with $A\in\Cc$. By Lemma~\ref{lem:nochain} a
non-minimal member of $\Cc$ is maximal, so $A\cup A'\in\Mm$, which is~(i).

Otherwise $|S_A\cap S_{A'}|=1$; by~\eqref{eq:pair} this forces $n$ odd and
equality throughout, i.e.\ $|S_A|=|S_{A'}|=\frac{n+1}{2}$ and
$S_A\cup S_{A'}=[n]$. Thus $|A|=|A'|=\frac{n-1}{2}$ and
$|A\cup A'|=n-1$. A set of size $n-1$ belonging to $\A$ is either a coatom or
would be a core set containing all of $T$; the latter is impossible by
Lemma~\ref{lem:under}. Hence $A\cup A'\in\Yy$. Finally
$|A|+|A'|=n-1=|A\cup A'|$ gives $A\cap A'=\emptyset$. This is~(ii).
\end{proof}

\begin{lemma}[Lines]\label{lem:lines}
Let $M\in\Mm$ and let $A\ne A'$ be members of $V_M$. Then $A\cup A'=M$.
Consequently $|V_M\cap V_{M'}|\le1$ for distinct $M,M'\in\Mm$, and every
internal pair is contained in $V_M$ for exactly one $M\in\Mm$, namely
$M=A\cup A'$.
\end{lemma}

\begin{proof}
$A\cup A'\subseteq M$ gives $|A\cup A'|\le|M|\le n-2$, so by
Lemma~\ref{lem:pair} the pair is internal and $A\cup A'\in\Mm$. Being contained
in the maximal set $M$, it equals $M$. If $A\ne A'$ both lay in
$V_M\cap V_{M'}$, we would get $M=A\cup A'=M'$.
\end{proof}

We call the sets $M\in\Mm$ with $k_M\ge2$ the \emph{lines}. By
Lemma~\ref{lem:lines} two distinct lines meet in at most one point of $\Cp$;
for $n$ even, Lemma~\ref{lem:pair} says moreover that every pair of points lies
on exactly one line, so that $(\Cp,\{V_M\})$ is a linear space.

\subsection*{Capacity of a line}

\begin{lemma}[Capacity]\label{lem:cap}
Let $M\in\Mm$. If $n$ is even then $k_M\le\lambda_M$. If $n$ is odd then
$k_M\le\lambda_M$, except when $\lambda_M=\frac{n-1}{2}$, in which case
$k_M\le\lambda_M+1$.
\end{lemma}

\begin{proof}
Write $\mu:=|M|=n-\lambda$ and, for $A\in V_M$ with $A\ne M$, let
$p_A:=\mu-|A|\ge1$ be its \emph{petal}. By Lemma~\ref{lem:children}(ii) the
petals of the members of $V_M\setminus\{M\}$ are pairwise disjoint nonempty
subsets of $M$, so their number is at most $\mu$ and $\sum p_A\le\mu$.

\emph{Case $\mu\le\nu$.} Then $M\in\Cp$, so $k_M\le\mu+1=n-\lambda+1$. Also
$\mu\le\nu$ gives $\lambda\ge n-\nu=\lfloor n/2\rfloor+1$, whence
$k_M\le n-\lambda+1\le\lceil n/2\rceil\le\lambda$ (for $n$ even the last chain
even gives $k_M\le n/2<\lambda$).

\emph{Case $\mu>\nu$.} Then $M\notin\Cp$, every $A\in V_M$ is a proper subset
of $M$ with $p_A\ge\mu-\nu\ge1$, and disjointness of the petals gives
$k_M(\mu-\nu)\le\mu$, i.e.
\begin{equation}\label{eq:cap}
  k_M\le\frac{n-\lambda}{n-\lambda-\nu}.
\end{equation}

Assume first $n$ even, so $\nu=n/2-1$ and $\mu\ge n/2$, i.e.\
$2\le\lambda\le n/2$. The bound~\eqref{eq:cap} is at most $\lambda$ if and only
if
\[
  n-\lambda\le\lambda(n-\lambda-\nu),
\]
i.e.
\[
  \lambda^2-\lambda\Bigl(2+\frac n2\Bigr)+n\le0,
\]
which is
\[
  (\lambda-2)\Bigl(\lambda-\frac n2\Bigr)\le0.
\]
This holds precisely on $2\le\lambda\le n/2$.

Now let $n$ be odd, so $\nu=(n-1)/2$ and $\mu\ge(n+1)/2$, i.e.\
$2\le\lambda\le(n-1)/2$. Put
\[
  u:=\mu-\nu=\frac{n+1}{2}-\lambda\ge1.
\]
Then~\eqref{eq:cap} reads
\[
  k_M\le\frac{\mu}{u}=1+\frac{(n-1)/2}{u}.
\]
If $u=1$, i.e.\ $\lambda=(n-1)/2$, this is
$k_M\le(n+1)/2=\lambda+1$. If $u\ge2$, it suffices, $k_M$ being an integer, to
show
\[
  1+\frac{n-1}{2u}<\lambda+1=\frac{n+3}{2}-u,
\]
i.e.
\[
  Q(u):=2u^2-u(n+1)+(n-1)<0.
\]
Here $2\le u\le(n-3)/2$. For $n=5$ this case cannot occur. For odd $n\ge7$,
$Q$ is convex with $Q(2)=5-n<0$ and $Q((n-3)/2)=5-n<0$; hence $Q<0$ on the
whole interval, and $k_M\le\lambda$.
\end{proof}

\subsection*{A global bound on \texorpdfstring{$q$}{q}}

\begin{lemma}[Fisher-type count]\label{lem:fisher}
\[
  \sum_{M\in\Mm}\binom{\lambda_M}{2}
  \le
  \binom n2-\binom m2
  =
  \binom t2+tm.
\]
\end{lemma}

\begin{proof}
By Lemma~\ref{lem:pairM} a pair $\{x,y\}\subseteq[n]$ is contained in $S_M$ for
at most one $M\in\Mm$, and by Lemma~\ref{lem:Zminus} no $S_M$ contains two
elements of $Z$, so no pair inside $Z$ is counted.
\end{proof}

\begin{lemma}[Global bound]\label{lem:q}
Let
\[
  f:=|\{M\in\Mm:\lambda_M=(n-1)/2\}|
\]
when $n$ is odd, and $f:=0$ when $n$ is even. Then
\[
  f\le\Bigl\lfloor\frac{4n}{n-3}\Bigr\rfloor,
\]
and $f=0$ unless $t\ge(n-3)/2$. Moreover
\[
  q(q-1)\le t(t-1)+2tm
  \qquad(n\text{ even}),
\]
and
\[
  q^2-2q\le t(t-1)+2tm+\frac32 f(n-1)
  \qquad(n\text{ odd}).
\]
\end{lemma}

\begin{proof}
For $f$: such sets $S_M$ have $(n-1)/2\ge2$ elements and pairwise meet in at
most one point, so
\[
  f\binom{(n-1)/2}{2}\le\binom n2
\]
by Lemma~\ref{lem:fisher}, which gives $f\le4n/(n-3)$; and
$\lambda_M\le t+1$ forces $(n-1)/2\le t+1$.

For $n$ even, every pair of distinct members of $\Cp$ is internal
(Lemma~\ref{lem:pair}) and lies on exactly one line (Lemma~\ref{lem:lines}), so
with Lemmas~\ref{lem:cap} and~\ref{lem:fisher}
\[
  \binom q2
  =
  \sum_{M\in\Mm}\binom{k_M}{2}
  \le
  \sum_{M\in\Mm}\binom{\lambda_M}{2}
  \le
  \binom t2+tm.
\]

For $n$ odd, let $X$ be the number of external pairs. By Lemma~\ref{lem:cap},
\[
  \sum_M\binom{k_M}{2}
  \le
  \sum_M\binom{\lambda_M}{2}+f\cdot\frac{n-1}{2},
\]
because $\binom{\lambda+1}{2}-\binom{\lambda}{2}=\lambda$.

For $X$, suppose $A\in\Cp$ has two external partners. By Lemma~\ref{lem:pair}
they are $S_A\setminus\{x\}$ and $S_A\setminus\{y\}$ for distinct $x,y$, and
their union is $S_A$, of size $(n+1)/2\le n-2$ for $n\ge5$. By
Lemma~\ref{lem:pair} again $S_A\in\Mm$, with $\lambda_{S_A}=(n-1)/2$. Distinct
such $A$ give distinct sets $S_A$, so at most $f$ members of $\Cp$ have external
degree at least $2$, and each external degree is at most
$k_{S_A}\le(n+1)/2$ by Lemma~\ref{lem:cap}. Hence
\[
  X\le\frac12\Bigl(q+f\frac{n-1}{2}\Bigr).
\]
Therefore
\[
  \binom q2
  =
  \sum_M\binom{k_M}{2}+X
  \le
  \binom t2+tm+f\frac{n-1}{2}+\frac q2+\frac{f(n-1)}4.
\]
Multiplying by $2$ and subtracting $q$ gives the stated odd bound.
\end{proof}

We write $Q(n,t)$ for the largest integer $q$ satisfying the relevant inequality
of Lemma~\ref{lem:q}; thus $q\le Q(n,t)$ always, and $Q(n,t)\le n$ for $n$ even.

\subsection*{The local bound}

\begin{lemma}[Local bound]\label{lem:local}
Assume $q\ge2$ and let $A_0\in\Cp$ be of maximum size. Put
\[
  a_0:=|A_0|,\qquad
  w_0:=n-2a_0,\qquad
  \sigma:=n-a_0,\qquad
  z_0:=|S_{A_0}\cap Z|.
\]
Thus $1\le w_0\le t+1$ and $z_0\ge m-a_0$. Let $\varepsilon$ be the number of
external partners of $A_0$. Then $\varepsilon=0$ unless $n$ is odd and
$w_0=1$, in which case $\varepsilon\le1$ if $f=0$ and
$\varepsilon\le(n+1)/2$ otherwise. Moreover:

\begin{enumerate}
\item[\rm(A)] If $A_0\in\Mm$ --- which forces $w_0\le2t+2-n$ --- then
\[
  \Phi^+\le w_0+\varepsilon+a_0(w_0+2).
\]

\item[\rm(B)] If $A_0\notin\Mm$, let $\mathcal L$ be the set of $M\in\Mm$ of the
form $M=A\cup A_0$ with $A\in\Cp\setminus\{A_0\}$, and let
$r:=|\mathcal L|$, so that $r+\varepsilon\le q-1$. Then every
$M\in\mathcal L$ satisfies $2\lambda_M\le n$ and $\lambda_M\le t+1$,
\[
  \sum_{M\in\mathcal L}\lambda_M(\lambda_M-1)
  \le
  \sigma(\sigma-1)-z_0(z_0-1)=:B,
\]
and
\[
  \Phi^+
  \le
  w_0+\varepsilon+\sum_{M\in\mathcal L}(2\lambda_M-w_0).
\]
Consequently
\[
  \Phi^+
  \le
  w_0+\varepsilon+
  \max_{0\le r\le q-1-\varepsilon}
  \Bigl[
    2\min\Bigl(
      \frac{r+\sqrt{r^2+4rB}}2,
      (t+1)r
    \Bigr)
    -rw_0
  \Bigr],
\]
the bracket being $0$ at $r=0$ (the case $\mathcal L=\varnothing$).
\end{enumerate}
\end{lemma}

\begin{proof}
The bound $w_0\le t+1$ follows by picking $A\in\Cp$, $A\ne A_0$: by
Lemma~\ref{lem:P},
\[
  2a_0\ge |A_0|+|A|\ge m-1,
\]
so $w_0=n-2a_0\le n-m+1=t+1$. The statement about $\varepsilon$ is contained in
the proof of Lemma~\ref{lem:q}: an external pair forces
$|A_0|=(n-1)/2$, i.e.\ $w_0=1$, and two external partners force
$S_{A_0}\in\Mm$ with $\lambda_{S_{A_0}}=(n-1)/2$, hence $f\ge1$; in that case
$\varepsilon\le k_{S_{A_0}}\le(n+1)/2$.

By Lemma~\ref{lem:pair} every $A\in\Cp\setminus\{A_0\}$ is either an external
partner of $A_0$ (and then $\wt(A)=1$) or satisfies $A\cup A_0\in\Mm$; in the
latter case $A\in V_M$ for $M=A\cup A_0$, and by Lemma~\ref{lem:lines} this $M$
is the unique line through $A_0$ containing $A$. Hence
\begin{equation}\label{eq:split-local}
  \Phi^+
  =
  w_0+\varepsilon+
  \sum_M\sum_{A\in V_M\setminus\{A_0\}}\wt(A),
\end{equation}
the outer sum being over the maximal core sets containing $A_0$ that carry a
further point.

Fix such an $M$ and write $\mu=|M|$. For $A\in V_M$ we have
\[
  \wt(A)=n-2|A|=n-2\mu+2p_A=2\lambda_M-n+2p_A,
\]
and since the petals are pairwise disjoint with $p_{A_0}=\mu-a_0$,
\begin{equation}\label{eq:petal}
  \sum_{A\in V_M\setminus\{A_0\}}p_A
  \le
  \mu-p_{A_0}
  =
  a_0.
\end{equation}

(A) Let $A_0\in\Mm$. Since $\Mm$ is an antichain, no $M\in\Mm$ properly contains
$A_0$, so the only maximal core set containing $A_0$ is $A_0$ itself: every
non-external $A\ne A_0$ satisfies $A\cup A_0=A_0$, i.e.\ $A\subsetneq A_0$.
Their number is at most $a_0$ by Lemma~\ref{lem:children}(ii), so by
\eqref{eq:split-local}, \eqref{eq:petal} and $2\lambda_{A_0}-n=w_0$,
\[
  \Phi^+\le w_0+\varepsilon+a_0w_0+2a_0.
\]
Finally $A_0\in\Mm$ gives $a_0\ge m-1$ by Lemma~\ref{lem:Zminus}, i.e.\
$w_0\le2t+2-n$.

(B) Let $A_0\notin\Mm$ and $M\in\mathcal L$. Then $A_0\subsetneq M$, and
$M\notin\Cp$: otherwise $M$ would be a member of $\Cp$ of size $>a_0$. Hence
$2\mu\ge n$, i.e.\ $2\lambda_M\le n$, and $k_M-1\ge1$, so by~\eqref{eq:petal}
\[
  \sum_{A\in V_M\setminus\{A_0\}}\wt(A)
  =
  (k_M-1)(2\lambda_M-n)+2\sum_{A\in V_M\setminus\{A_0\}}p_A
  \le
  (2\lambda_M-n)+2a_0
  =
  2\lambda_M-w_0.
\]
The sets $S_M$, $M\in\mathcal L$, are subsets of $S_{A_0}$ which pairwise meet
in at most one point (Lemma~\ref{lem:pairM}) and none of which contains two
points of $Z$ (Lemma~\ref{lem:Zminus}); double counting the pairs inside
$S_{A_0}$ gives
\[
  \sum_{M\in\mathcal L}\binom{\lambda_M}{2}
  \le
  \binom{\sigma}{2}-\binom{z_0}{2}.
\]
Distinct lines through $A_0$ carry disjoint sets of further points, and the
$\varepsilon$ external partners of $A_0$ lie on none of them, so
$r+\varepsilon\le q-1$. For fixed $r$, maximising $\sum\lambda_M$ subject to
$\sum\lambda_M(\lambda_M-1)\le B$ gives, by Cauchy--Schwarz,
\[
  \sum\lambda_M
  \le
  \frac12\bigl(r+\sqrt{r^2+4rB}\bigr),
\]
while $\lambda_M\le t+1$ gives $\sum\lambda_M\le(t+1)r$. Substituting into
\eqref{eq:split-local} yields the last display.
\end{proof}

\subsection*{Closing the window}

\begin{definition}\label{def:psi}
For integers $n\ge4$ and $1\le t\le n$ put $m=n-t$ and let
$R:=Q(n,t)-1$. For each integer $w_0$ with $1\le w_0\le t+1$ and
$w_0\equiv n\pmod2$, set
\[
  a_0=\frac{n-w_0}{2},\qquad
  \sigma=\frac{n+w_0}{2},\qquad
  z_0=\max(0,m-a_0),
\]
\[
  B=\sigma(\sigma-1)-z_0(z_0-1),
\]
and
\[
  \varepsilon(w_0)=
  \begin{cases}
    0,& n\text{ even or }w_0\ge2,\\
    1,& n\text{ odd},\ w_0=1,\ f_{\max}(n,t)=0,\\
    \frac{n+1}{2},&\text{otherwise},
  \end{cases}
\]
where
\[
  f_{\max}(n,t)=
  \begin{cases}
    \left\lfloor\frac{4n}{n-3}\right\rfloor,
      & n\text{ odd and }t+1\ge\frac{n-1}{2},\\
    0,&\text{otherwise}.
  \end{cases}
\]
Define
\[
  \Psi_A(w_0)=
  \begin{cases}
    w_0+\varepsilon(w_0)+a_0(w_0+2),
      & w_0\le2t+2-n,\\
    -\infty,&\text{otherwise},
  \end{cases}
\]
and
\[
  \Psi_B(w_0)=
  w_0+\varepsilon(w_0)+
  \max_{0\le r\le\max(0,R-\varepsilon(w_0))}
  \Bigl[
    2\min\Bigl(
      \frac{r+\sqrt{r^2+4rB}}2,
      (t+1)r
    \Bigr)
    -rw_0
  \Bigr],
\]
the bracket being $0$ at $r=0$. Finally put
\[
  \Psi(n,t):=\max_{w_0}\max\{\Psi_A(w_0),\Psi_B(w_0)\}.
\]
\end{definition}

\begin{remark}\label{rem:r0}
The value $r=0$ has to be admitted in $\Psi_B$, because $\mathcal L$ may be
empty; it never dominates. Indeed, if $r=0$ then every member of
$\Cp\setminus\{A_0\}$ is an external partner of $A_0$, so $q\ge2$ forces
$\varepsilon\ge1$, hence $n$ is odd and $w_0=1$ by Lemma~\ref{lem:local}. In
that case $R\ge q-1\ge\varepsilon+0$, so $r=1$ is admissible whenever
$R-\varepsilon\ge1$, and the bracket at $r=1$ equals
\[
  2\min\Bigl(\frac{1+\sqrt{1+4B}}2,\,t+1\Bigr)-1\ \ge\ 2\cdot 1-1\ =\ 1\ >\ 0,
\]
since $B\ge0$ and $t\ge1$. Thus admitting $r=0$ changes $\Psi$ only in the
degenerate case $R=\varepsilon$, where it is needed.
\end{remark}

\begin{proposition}\label{prop:psi}
For every separating union-closed family of height four with $4\le n$ on $[n]$,
\[
  \Phi^+\le\max\{n,\Psi(n,t)\}.
\]
\end{proposition}

\begin{proof}
If $q\le1$ then $\Phi^+\le n$. If $q\ge2$, apply Lemma~\ref{lem:local} and
$q\le Q(n,t)$ from Lemma~\ref{lem:q}.
\end{proof}

\begin{theorem}\label{thm:window}
For every pair $(n,t)$ with $n\ge11$, $t\ge3$ and $n\le4t+3$,
\[
  \Psi(n,t)<t(n-2)+n.
\]
\end{theorem}

\begin{proof}
For $11\le n\le200$ this is a finite evaluation of Definition~\ref{def:psi}; it
is carried out by the script \texttt{verify\_h4.py} in the repository, using
exact integer upper bounds for the square-root term (so that the verification
involves no floating-point arithmetic). The largest ratio
$\Psi(n,t)/(t(n-2)+n)$ over this range is attained at $(n,t)=(11,4)$, where it
is $<0.73$; it decreases steadily with $n$, being $<0.30$ already at $n=89$.

For $n\ge201$ we argue directly. The case $r=0$ of $\Psi_B$ contributes only
$w_0+\varepsilon\le(n+1)+\frac{n+1}2$, which is below the target for $n\ge5$,
so we may assume $r\ge1$. In case~(B), $w_0\ge1$ gives
\[
  2\min\Bigl(
    \frac{r+\sqrt{r^2+4rB}}2,
    (t+1)r
  \Bigr)-rw_0
  \le
  \sqrt{r^2+4rB}
  \le
  r+2\sqrt{rB}.
\]
Also $B\le\sigma(\sigma-1)<\sigma^2\le n^2$, since $\sigma=n-a_0\le n$.

Next we bound $r$. For $n$ even, $t(t-1)+2tm=t(2n-t-1)$ is maximised over
$0\le t\le n$ at $t\in\{n-1,n\}$, with value $n^2-n$; hence
$q(q-1)\le n^2-n$ and $q\le n$. For $n$ odd we have
$f\le\lfloor 4n/(n-3)\rfloor\le5$ for $n\ge11$, so Lemma~\ref{lem:q} gives
\[
  (q-1)^2\le q^2-2q+1\le (n^2-n)+\tfrac{15}2(n-1)+1
  = n^2+\tfrac{13}2 n-\tfrac{13}2
  <\Bigl(n+\tfrac{13}4\Bigr)^2 ,
\]
so $q<n+\frac{17}4$ and $q\le n+4$ in either case. Since $r\le q-1-\varepsilon$
we get $r\le n+3$. Therefore, using $w_0\le t+1\le n+1$ and
$\varepsilon\le(n+1)/2$,
\[
  \Psi_B(w_0)
  \le
  (n+1)+\frac{n+1}{2}+(n+3)+2n\sqrt{n+3}
  \le
  \frac52n+\frac{11}{2}+2n\sqrt{n+4}.
\]
For case~(A), $a_0=\frac{n-w_0}2$ gives
$a_0(w_0+2)=\frac{(n-w_0)(w_0+2)}2\le\frac{(n+2)^2}8$, the maximum of the
quadratic being at $w_0=\frac{n-2}2$; hence
\[
  \Psi_A(w_0)\le\frac{3(n+1)}2+\frac{(n+2)^2}{8}.
\]

On the other hand $t\ge(n-3)/4$ gives
\[
  t(n-2)+n
  \ge
  \frac{(n-3)(n-2)}4+n
  =
  \frac{n^2-n+6}{4}.
\]

For case~(A) it therefore suffices that
$\frac{(n+2)^2}8+\frac{3(n+1)}2\le\frac{n^2-n+6}4$; multiplying by $8$, this
reads $n^2+16n+16\le 2n^2-2n+12$, i.e.\ $n^2-18n-4\ge0$, which holds for
$n\ge19$.

For case~(B) it suffices that
$\frac52n+\frac{11}2+2n\sqrt{n+4}\le\frac{n^2-n+6}4$, i.e.\ that
\[
  8n\sqrt{n+4}\le n^2-11n-16,
  \qquad\text{equivalently}\qquad
  64n^3+256n^2\le\bigl(n^2-11n-16\bigr)^2 ,
\]
both sides of the first inequality being nonnegative for $n\ge13$. Now
$n\ge192$ gives $11n+16\le12n\le\frac{n^2}{16}$, whence
$n^2-11n-16\ge\frac{15}{16}n^2$ and
$\bigl(n^2-11n-16\bigr)^2\ge\frac{225}{256}n^4$; and $n\ge128$ gives
$256n^2\le2n^3$, whence $64n^3+256n^2\le66n^3$. Since
$\frac{225}{256}n^4\ge66n^3$ for $n\ge\frac{66\cdot256}{225}=75.1\ldots$, all
three conditions hold for $n\ge201$, which proves the claim in that range.

The ranges $11\le n\le200$ and $n\ge201$ are complementary, so the theorem
holds for every $n\ge11$.
\end{proof}

\begin{corollary}\label{cor:residual}
For every separating union-closed family of height four with $4\le n$, the
residual inequality of Theorem~\ref{thm:residual} holds.
\end{corollary}

\begin{proof}
By Theorem~\ref{thm:small} we may assume $n\ge11$ and exclude the cases
$n=11$, $t\le3$. If $n\ge4t+4$, we use Lemma~\ref{lem:q} to bound $q$. Since
$t\ge1$,
we have $n\ge8$ and $t\le(n-4)/4$. For odd $n$, Lemma~\ref{lem:q} gives
$f=0$: indeed $t\le(n-4)/4<(n-3)/2$, below the threshold required for $f>0$.
In either parity, the right-hand side in the bound of Lemma~\ref{lem:q} is
therefore at most
\[
  t(t-1)+2tm=t(2n-t-1)
  \le \frac{(n-4)(2n-1)}4
  <(n-1)(n-3),
\]
where the strict inequality follows from $2n^2-7n+8>0$. For even $n$, if
$q\ge n-1$ then $q(q-1)\ge(n-1)(n-2)>(n-1)(n-3)$, contradicting
Lemma~\ref{lem:q}. For odd $n$, $q\ge n-1$ would give
$q^2-2q\ge(n-1)(n-3)$, again a contradiction. Thus $q\le n-2$.

If $q=0$ then $\Phi^+=0$. If $q=1$ then $\Phi^+\le n\le t(n-2)+n$.
If $q\ge2$, Theorem~\ref{thm:master} gives
\[
  \Phi^+\le q(t+1)\le(n-2)(t+1)=t(n-2)+n-2.
\]
Otherwise $n\le4t+3$. Since $n\ge11$, $t=1$ is impossible, and $t=2$ would
force $n=11$, a case already excluded by Theorem~\ref{thm:small}; hence $t\ge3$.
Then
Theorem~\ref{thm:window} together with Proposition~\ref{prop:psi} gives
\[
  \Phi^+\le\max\{n,\Psi(n,t)\}<t(n-2)+n,
\]
because $t\ge3$ and $n\ge11$ imply $n<t(n-2)+n$.
\end{proof}

\subsection*{Why the counting bound $q\le n-2$ is not enough}

The reduction via Corollary~\ref{cor:count} is useful but cannot be the final
route: the sufficient condition $q\le n-2$ is false inside the residual window.

\begin{proposition}\label{prop:qcex}
Let $n\ge6$ be even, let $[n]=M_1\sqcup M_2$ with
$|M_1|=|M_2|=n/2$, and let $\A$ be the union closure of
\[
  \{M_1\setminus\{x\}:x\in M_1\}
  \cup
  \{M_2\setminus\{y\}:y\in M_2\}.
\]
Then $\A$ is separating and union-closed with $h(\A)=4$, $t=n$, $q=n$, and
$\Phi^+=2n$.
\end{proposition}

\begin{proof}
The closure consists of the $n$ generators, the two sets $M_i$, the sets
$[n]\setminus\{x,y\}$ with $x\in M_1$, $y\in M_2$, all $n$ coatoms, and $[n]$.
Union-closedness and separation are immediate, and a maximal chain
\[
  M_1\setminus\{x\}
  \subsetneq
  M_1
  \subsetneq
  [n]\setminus\{y\}
  \subsetneq
  [n]
\]
has four members while $\Cc$ contains no three-element chain. The members of
size below $n/2$ are exactly the $n$ generators, each of weight $2$.
\end{proof}

For $n\ge12$ these families satisfy $n\ge11$, $t\ge3$, $n\le4t+3$, and
$|A\cup A'|\ge(n-1)/2$ for all distinct $A,A'\in\Cp$, i.e.\ they lie in the
residual window, and they have $q=n>n-2$. They also show that
Lemma~\ref{lem:q} is sharp, since
$q(q-1)=n(n-1)=t(t-1)+2tm$ there; the second family attaining equality in
Lemma~\ref{lem:q} is that of Example~\ref{ex:sharp}, where
$q(q-1)=6=t(t-1)+2tm$.

\section{The height-four theorem and concluding remarks}\label{sec:main}

\begin{theorem}\label{thm:main}
Let $\A\subseteq2^{[n]}$ be separating and union-closed with $h(\A)=4\le n$. Then
\[
  \Avg(\A)\ge\frac n2.
\]
In particular, Frankl's conjecture holds for every separating union-closed
family of height four with $4\le n$.
\end{theorem}

\begin{proof}
By Corollary~\ref{cor:residual} we have $\Phi^+\le t(n-2)+n$. Since
$\Phi\le\Phi^+$, Proposition~\ref{prop:reform} gives $\Avg(\A)\ge n/2$.
Frankl's conjecture for this class follows from the standard double-counting
identity: if $\Avg(\A)\ge n/2$, then some element lies in at least $|\A|/2$
member sets.
\end{proof}

\begin{theorem}\label{thm:t2-final}
Every separating union-closed family of height four with $4\le n$ and $t\le2$ satisfies
$\Avg(\A)\ge n/2$.
\end{theorem}

\begin{proof}
This follows from Corollary~\ref{cor:residual}, which proves the stronger
residual inequality throughout the stated range.
\end{proof}

\begin{remark}
The proof bounds $\Phi^+$ line by line rather than through the cardinality
$q=|\Cp|$. This is unavoidable: Corollary~\ref{cor:count} gives a convenient
sufficient condition, but Proposition~\ref{prop:qcex} shows that $q\le n-2$ is
not available inside the residual window.
\end{remark}

\begin{remark}
Earlier drafts formulated the remaining problem as a per-coatom bound: by
Proposition~\ref{prop:identity}, $\Phi^+\le t(n-2)+n$ would follow from
\[
  \widetilde C_x\le(n-2)+\frac nt
  \qquad(x\in T).
\]
That reformulation remains useful, but the closing argument above does not need
it. The pair-line structure of $\Cp$ gives a direct global bound on $\Phi^+$.
\end{remark}

\begin{remark}
Since $h=4$ is the last height at which averaging can work
\cite[Thm.~3.2]{Bouchard}, Theorem~\ref{thm:main} is the natural end point of
the averaging method for separating union-closed families.
\end{remark}

\subsection*{Machine verification}

Every statement of Section~\ref{sec:closing} has been tested independently of
its proof. The script \texttt{verify\_h4.py} enumerates separating union-closed
families of height four from the definitions alone --- union-closedness,
separation, and a longest chain of exactly four member sets --- and then checks
Lemmas~\ref{lem:pair}, \ref{lem:lines}, \ref{lem:cap}, \ref{lem:fisher},
\ref{lem:q} and~\ref{lem:local}, the assembled bound of
Lemma~\ref{lem:local}, the value $\Psi(n,t)$ of Definition~\ref{def:psi}, the
residual inequality of Theorem~\ref{thm:residual} and the conclusion
$\Avg(\A)\ge n/2$ as assertions. The same assertions were added to the
C enumerator used for Computation~I. No violation was found in any of the
following.

\begin{itemize}
\item The complete enumerations for $n=4$ ($2\,034$ families), $n=5$
      ($346\,028$ families) and $n=6$ with $t\le5$ ($106\,040\,636$ families),
      whose counts agree with those of Section~\ref{sec:conj}.
\item The families of Proposition~\ref{prop:qcex} for $8\le n\le18$, which lie
      inside the residual window and have $q=n>n-2$, together with several
      hundred further families constructed inside that window with
      $10\le q\le14$.
\item Random samples for $6\le n\le16$.
\end{itemize}

The window evaluation of Theorem~\ref{thm:window} uses exact integer upper
bounds for the square-root term, so no floating-point arithmetic enters the
verification; the candidate-point maximisation over $r$ was checked against a
full integer scan for all $(n,t,w_0)$ with $n\le120$.

\subsection*{Use of AI assistance}
Large language models (Qwen, Alibaba Group; Claude, Anthropic) were used as a
research assistant during the preparation of this work: for exploring proof
strategies, for drafting and checking the verification code, and for language
editing and \LaTeX\ preparation. The models were not used as authors and are not
listed as such, in accordance with arXiv policy. All definitions, statements,
proofs and computational verifications were formulated and validated by the
author, who takes full responsibility for the correctness of the paper.

\subsection*{Data availability}
All code, verification logs and the archived version of this paper are
publicly available at Zenodo under DOI
\url{https://doi.org/10.5281/zenodo.21739198}.
The development repository is
\url{https://github.com/YSKpython/frankl-height4-avg-bound}.

\begin{thebibliography}{99}

\bibitem{AHS} R.~Alweiss, B.~Huang and M.~Sellke,
\emph{Improved lower bound for Frankl's union-closed sets conjecture},
Electron. J. Combin. \textbf{31} (2024), \#P3.35.

\bibitem{BBE} I.~Balla, B.~Bollob\'as and T.~Eccles,
\emph{Union-closed families of sets},
J. Combin. Theory Ser. A \textbf{120} (2013), 531--544.

\bibitem{BM} I.~Bo\v{s}njak and P.~Markovi\'c,
\emph{The $11$-element case of Frankl's conjecture},
Electron. J. Combin. \textbf{15} (2008), \#R88.

\bibitem{Bouchard} C.~Bouchard,
\emph{An averaging result for union-closed families of sets},
preprint (2025), \texttt{arXiv:2509.12537}.

\bibitem{BruhnSchaudt} H.~Bruhn and O.~Schaudt,
\emph{The journey of the union-closed sets conjecture},
Graphs Combin. \textbf{31} (2015), 2043--2074.

\bibitem{Cambie} S.~Cambie,
\emph{Progress on the union-closed conjecture and offsprings in winter
2022--2023}, preprint (2023), \texttt{arXiv:2306.12351}.

\bibitem{Czedli09} G.~Cz\'edli,
\emph{On averaging Frankl's conjecture for large union-closed sets},
J. Combin. Theory Ser. A \textbf{116} (2009), 724--729.

\bibitem{CzedliMS} G.~Cz\'edli, M.~Mar\'oti and E.~T.~Schmidt,
\emph{On the scope of averaging for Frankl's conjecture},
Order \textbf{26} (2009), 31--48.

\bibitem{FalgasRavry} V.~Falgas-Ravry,
\emph{Minimal weight in union-closed families},
Electron. J. Combin. \textbf{18} (2011), \#P95.

\bibitem{Frankl95} P.~Frankl,
\emph{Extremal set systems}, in: R.~L.~Graham, M.~Gr\"otschel and L.~Lov\'asz
(eds.), Handbook of Combinatorics, vol.~2, Elsevier, 1995, pp.~1293--1329.

\bibitem{Gilmer} J.~Gilmer,
\emph{A constant lower bound for the union-closed sets conjecture},
preprint (2022), \texttt{arXiv:2211.09055}.

\bibitem{HachimoriKashiwabara} M.~Hachimori and K.~Kashiwabara,
\emph{On the averaging problem of ideal families related to Frankl's conjecture
with formal proof by Lean 4}, preprint (2025), \texttt{arXiv:2504.13454}.

\bibitem{Poonen} B.~Poonen,
\emph{Union-closed families},
J. Combin. Theory Ser. A \textbf{59} (1992), 253--268.

\bibitem{Reimer} D.~Reimer,
\emph{An average set size theorem},
Combin. Probab. Comput. \textbf{12} (2003), 89--93.

\bibitem{Sawin} W.~Sawin,
\emph{An improved lower bound for the union-closed set conjecture},
preprint (2023), \texttt{arXiv:2211.11504}.

\bibitem{Studer} L.~Studer,
\emph{An asymptotic version of the union-closed sets conjecture},
preprint (2020), \texttt{arXiv:2005.00361}.

\bibitem{Tian} C.~Tian,
\emph{Union-closed sets conjecture holds for height no more than $3$ and height
no less than $N-1$}, preprint (2022), \texttt{arXiv:2112.06659}.

\end{thebibliography}

\end{document}
